\documentclass[12pt,a4paper]{article}
\usepackage[utf8]{inputenc}
\usepackage[english]{babel}
\usepackage{geometry}
\geometry{margin=2.5cm}
\usepackage{hyperref}
\usepackage{graphicx}
\usepackage{amsmath}
\usepackage{amssymb}
\usepackage{booktabs}
\usepackage{enumitem}
\usepackage{natbib}
\usepackage{algorithm}
\usepackage{algpseudocode}

\title{Reward Shape and Frame Stacking in DQN:\\A Reward $\times$ Observation Ablation in Multi-Objective Grid Survival}
\author{Philipp Merz}
\date{2026}

\begin{document}
\maketitle
\begin{abstract}
Multi-objective embodied control tasks such as legged-robot locomotion and aerial vehicle navigation typically rely on dense per-step reward shaping with multiple weighted components, whose tuning is a known practical bottleneck. We study how reward shape and short-horizon temporal context interact in a DQN agent on a multi-objective grid-survival task. A $3 \times 2 \times 4$ ablation crosses three reward configurations (\textsc{baseline}, \textsc{absolute\_proximity}, \textsc{minimal}) with two observation types (single-frame and 4-frame stacked), four seeds per cell. Three findings emerge. First, frame stacking has a reward-dependent direction: it adds $71$ ticks of mean survival to \textsc{minimal} but degrades \textsc{absolute\_proximity} and \textsc{baseline} by $42$ and $91$ ticks respectively. Second, without frame stacking, performance is largely independent of reward shape and clusters at $634$--$674$ mean survival. Third, the closest cell, \textsc{minimal} with frame stacking, reaches statistical parity with a hand-tuned scripted reference in three of four seeds and significantly loses in the fourth, but does not exceed it. We interpret the asymmetry as preliminary evidence that frame stacking is not an unconditional remedy for under-performance; under our experimental conditions its benefit is contingent on the reward function having clean gradient structure. We frame the result as a single-environment characterization that motivates further controlled comparisons rather than a general claim.
\end{abstract}

\tableofcontents
\newpage

% =============================================================================
\section{Introduction}
\label{sec:introduction}
% =============================================================================

Reinforcement learning is increasingly applied to embodied control tasks where success depends on balancing multiple competing objectives. Legged-robot locomotion policies must trade forward velocity against energy use, joint-limit avoidance, foot-clearance, and stability \citep{Hwangbo2019}; quadrotor and unmanned aerial vehicle controllers balance trajectory tracking against control smoothness, motor saturation, and obstacle avoidance \citep{Hwangbo2017, Loquercio2021}. These tasks are \emph{multi-objective}: a single episode is judged by a combination of progress, safety, and resource use, with no single scalar return that cleanly summarises performance. They are also sparsely rewarded: the natural environment-driven reward (robot fell, drone crashed, time-out) arrives only at episode end, creating a so-called \emph{credit-assignment problem} because many low-level actions separate the start of an episode from the rewarded terminal state. The standard remedy is to augment the terminal reward with dense per-step shaping rewards intended to steer behaviour toward a specific balance of competing goals. For example, the published reward functions for the works above each combine six to twelve weighted terms, and tuning these weights is widely reported as one of the dominant practical costs of training such policies.

Dense shaping introduces its own design problem. The relative magnitudes between components determine what the agent actually optimises, and small changes can produce qualitatively different behaviours. Well-known failure modes called \emph{reward hacking} \citep{Amodei2016} emerge when the learned policy maximizes these dense rewards in unanticipated and unintended ways.

A natural complement to this reward design is the observation. Many environments expose only the current state, which carries position information but not motion. From a single frame, the agent cannot tell an approaching threat from a retreating one. Frame stacking \citep{Mnih2015}, done by concatenating the last $k$ frames so the network can infer derived quantities like velocity directly, serves as the standard solution to this problem. However, the interaction between reward shape and frame stacking is not well understood.

\paragraph{Scope.} We isolate the reward-design problem in a tractable gridworld testbed. According to \citet{Leike2017} and \citet{ChevalierBoisvert2023} gridworlds provide contexts for studying RL phenomena that are harder to disentangle in continuous-action embodied settings. The task is a multi-objective survival environment where a DQN agent must eat, evade enemies, and use shelters; the comparison target is a scripted heuristic baseline. We run a $3 \times 2 \times 4$ ablation crossing three reward configurations (\textsc{baseline}, \textsc{absolute\_proximity}, \textsc{minimal}; defined in Section~\ref{sec:reward_configs}) with two observation types (single-frame and 4-frame stacked), four random seeds per cell. We add one illustrative single-seed run, \textsc{weak\_proximity}, and a four-seed, $50{,}000$-tick-per-episode training of the headline cell as a robustness check on the choice of episode cap.

\paragraph{Findings.} \textsc{minimal} with frame stacking is the highest-performing of the six cells but does not exceed the heuristic: it reaches statistical parity in three of four seeds and significantly loses in the fourth, with mean survival $745$ ticks against the heuristic's $768$. Frame stacking does not improve every reward configuration uniformly: with \textsc{minimal} it improves mean survival by $71$ ticks; with \textsc{absolute\_proximity} it degrades survival by $42$ ticks; with \textsc{baseline} it degrades survival by $91$ ticks. Adding observation history is therefore not an unconditional improvement under our experimental conditions; under leaky reward designs it actively hurts. We discuss plausible mechanisms in Section~\ref{sec:discussion}.

\paragraph{Agent and compute.} The agent is a Rainbow-lite DQN \citep{Hessel2018} combining Double DQN, a Dueling architecture, prioritized experience replay, and 3-step returns, along with stability techniques (LayerNorm, AdamW, Polyak target updates). Training runs on single consumer hardware, following work by \citet{ObandoCeron2021}, \citet{Schwarzer2023}, \citet{ObandoCeron2024}, and \citet{Clark2024} demonstrating that algorithmic choices can substitute for large-scale compute. The architecture, all hyperparameters except the frame-stack size, and the training budget are held constant across all twenty-four main-ablation cells.

Section~\ref{sec:background} reviews the relevant background. Section~\ref{sec:environment} describes the survival environment. Section~\ref{sec:architecture} details the agent architecture. Section~\ref{sec:experiments} presents the experiments, including failure-mode analysis, reward-design heuristics, and the main ablation. Sections~\ref{sec:discussion} and~\ref{sec:conclusion} discuss findings and conclude.

% =============================================================================
\section{Background}
\label{sec:background}
% =============================================================================

\subsection{Reinforcement Learning Fundamentals}

We formulate the task as a Markov Decision Process (MDP) defined by a tuple $(S, A, P, R, \gamma)$: state space $S$, action space $A$, transition function $P(s' \mid s, a)$, reward function $R(s, a, s')$, and discount factor $\gamma \in [0, 1)$ \citep{SuttonBarto2018}.

The agent interacts with the environment in discrete time steps $t = 0, 1, 2, \ldots$, receiving reward $r_t \in \mathbb{R}$ at each step. The goal is to find a policy $\pi: S \to A$ that maximizes the expected discounted return $G_t = \sum_{k=0}^{\infty} \gamma^k r_{t+k+1}$. The action-value function
\begin{equation}
Q^\pi(s, a) = \mathbb{E}_\pi\!\left[G_t \mid s_t = s,\; a_t = a\right]
\end{equation}
gives the expected return of taking action $a$ in state $s$ and following $\pi$ thereafter. The optimal action-value function $Q^*(s, a) = \max_\pi Q^\pi(s, a)$ satisfies the Bellman optimality equation:
\begin{equation}
Q^*(s, a) = \mathbb{E}\!\left[r_t + \gamma \max_{a'\in A} Q^*(s_{t+1}, a')\right],
\end{equation}
where $s_{t+1}$ is the successor state reached after taking $a$ in $s$.

Q-learning estimates $Q^*$ directly through temporal difference (TD) updates:
\begin{equation}
Q(s_t, a_t) \leftarrow Q(s_t, a_t) + \alpha\!\left[r_t + \gamma \max_{a'} Q(s_{t+1}, a') - Q(s_t, a_t)\right],
\end{equation}
where $\alpha$ is the learning rate. The bracketed quantity is the TD error $\delta_t$.

\subsection{Deep Q-Networks}

When the state space is large, tabular Q-learning is infeasible. Deep Q-Networks (DQN) approximate $Q^*$ with a neural network $Q(s, a; \theta)$ parameterized by weights $\theta$ \citep{Mnih2013, Mnih2015}.

Two mechanisms stabilize training. First, \emph{experience replay} stores transitions $(s, a, r, s')$ in a buffer and samples mini-batches uniformly, breaking temporal correlations. Second, a \emph{target network} with parameters $\theta^-$ provides stable TD targets. The loss for a sampled transition is:
\begin{equation}
L(\theta) = \left[r + \gamma \max_{a'} Q(s', a'; \theta^-) - Q(s, a; \theta)\right]^2.
\end{equation}
The target parameters $\theta^-$ are periodically copied from $\theta$.

Exploration uses an $\varepsilon$-greedy policy: with probability $\varepsilon$ a random action is taken, otherwise the greedy action $\arg\max_a Q(s, a; \theta)$.

\subsection{DQN Extensions}
\label{sec:dqn_extensions}

Several extensions address specific weaknesses of the basic algorithm.

\paragraph{Double DQN.} Standard DQN overestimates action values because the same network both selects and evaluates the maximizing action. Double DQN \citep{vanHasselt2016} decouples these steps:
\begin{equation}
y = r + \gamma\, Q\!\left(s',\; \arg\max_{a'} Q(s', a'; \theta);\; \theta^-\right).
\end{equation}
The online network selects the action; the target network evaluates it.

\paragraph{Dueling architecture.} The Dueling network \citep{Wang2016} decomposes the Q-function into a state value $V(s)$ and an action advantage $A(s, a)$:
\begin{equation}
Q(s, a) = V(s) + A(s, a) - \frac{1}{|A|}\sum_{a'} A(s, a'),
\end{equation}
where $|A|$ is the number of actions. This decomposition allows the network to learn state values without evaluating every action, which helps in states where the action choice matters little.

\paragraph{Prioritized experience replay.} Uniform replay treats all transitions equally. Prioritized experience replay (PER) \citep{Schaul2016} samples transitions proportionally to their TD error magnitude, revisiting surprising transitions more often. The sampling probability of transition $i$ is:
\begin{equation}
P(i) = \frac{p_i^\alpha}{\sum_k p_k^\alpha}, \quad p_i = |\delta_i| + \epsilon_p,
\end{equation}
where $\alpha$ controls prioritization strength and $\epsilon_p$ is a small constant ensuring nonzero probability. Non-uniform sampling introduces bias, corrected by importance-sampling (IS) weights:
\begin{equation}
w_i = \left(\frac{1}{N \cdot P(i)}\right)^\beta\!\bigg/\; \max_j w_j,
\end{equation}
where $N$ is the number of stored transitions and $\beta$ is annealed from a starting value to 1 during training, gradually applying the full bias correction.

\paragraph{Multi-step returns.} Instead of bootstrapping from the immediate next state, $n$-step returns \citep{Sutton1988} accumulate rewards over $n$ steps:
\begin{equation}
\label{eq:nstep}
R_t^{(n)} = \sum_{k=0}^{n-1} \gamma^k r_{t+k+1} + \gamma^n \max_{a'} Q(s_{t+n}, a'; \theta^-).
\end{equation}
This propagates reward information faster at the cost of increased variance. Moderate values of $n$ (typically 2--5) improve learning speed in practice.

\paragraph{Rainbow.} \citet{Hessel2018} combined six DQN extensions into a single agent: Double DQN, Dueling, PER, multi-step returns, distributional RL (C51), and NoisyNet. An ablation study showed that each component contributes, with prioritized replay and multi-step returns providing the largest individual gains.

We adopt four of these six components: Double DQN, Dueling, PER, and 3-step returns. We omit distributional RL due to implementation complexity and NoisyNet because $\varepsilon$-greedy exploration is sufficient at smaller scales \citep{ObandoCeron2021, Fortunato2018}.

\subsection{Training Stability Techniques}
\label{sec:stability}

\paragraph{Polyak soft target updates.} The original DQN copies $\theta$ to $\theta^-$ every $C$ steps, creating discontinuities in the targets. Polyak averaging \citep{Polyak1992, Lillicrap2016} replaces the hard copy with an exponential moving average, updating the target at every step:
\begin{equation}
\theta^- \leftarrow (1 - \tau)\,\theta^- + \tau\,\theta,
\end{equation}
with a small $\tau$. This produces smoother target evolution and removes the need to tune $C$.

\paragraph{AdamW.} The Adam optimizer applies L2 regularization to the gradient, which interacts poorly with adaptive learning rates. AdamW \citep{Loshchilov2019} decouples weight decay from the gradient update, providing effective regularization that helps maintain network plasticity during long training runs.

\paragraph{LayerNorm.} Batch normalization is unreliable in RL because the data distribution is non-stationary and effective batch sizes are small. LayerNorm \citep{Ba2016} normalizes across features within each sample, independent of the batch, making it suitable for the single-sample updates common in RL.

\subsection{Reward Shaping}
\label{sec:reward_shaping_bg}

When multiple reward components are combined with different weights, the agent maximizes the weighted sum, which may not correspond to any single intended behavior. This kind of mismatch is commonly described as reward hacking \citep{Amodei2016}. In a survival setting, it creates a tension: a high food reward produces aggressive foraging at the cost of safety, a high death penalty produces shelter-seeking at the cost of starvation, and a high shelter reward produces passive camping. The interaction between these components, not just their individual magnitudes, determines the learned policy.

This interaction is the central object of study. We fix the agent and vary only the reward weights to isolate their effect on behavior.

% =============================================================================
\section{Environment}
\label{sec:environment}
% =============================================================================

Grid-based environments are a common testbed for reinforcement learning, from safety-focused gridworlds \citep{Leike2017} to the modular Minigrid framework \citep{ChevalierBoisvert2023}. Our environment extends this tradition with a multi-objective survival task that foregrounds the reward design problem.

\subsection{World Design}
\label{sec:world_design}

The environment is a $64 \times 64$ bounded grid with discrete time steps. An episode ends when the agent dies (health reaches zero) or when an episode cap on the number of ticks is reached. The choice of cap and its empirical non-bindingness are discussed in Section~\ref{sec:setup}. The world contains four entity types.

\paragraph{Agent.} A single agent with two resource bars: health $h$ and hunger $u$, both with maximum value $h_{\max}=u_{\max}=100$. The agent spawns at the grid center with both bars full. Hunger depletes by 0.2 per tick and by an additional 0.5 when moving. When hunger reaches zero, the agent starves and loses 0.5 health per tick. When hunger exceeds 30\% of maximum, health regenerates at 0.5 per tick up to the cap.

\paragraph{Enemies.} Mobile entities that deal 15 damage on contact (Manhattan distance $\leq 1$). Each tick, an enemy moves with probability 0.5. Within a vision range of 3 (Manhattan distance), enemies pursue the agent along the axis of greatest distance. Outside this range, they move randomly. Enemies cannot enter shelter tiles or tiles occupied by another enemy. The world supports up to 8 enemies simultaneously, with roughly a third of this capacity present at episode start. New enemies appear with probability 0.01 per tick.

\paragraph{Food.} Static items that restore 20 hunger when consumed by stepping onto their tile. Up to 36 food items exist simultaneously, with half present at episode start. New food appears with probability 0.05 per tick.

\paragraph{Shelters.} Static tiles generated at episode start with density 0.02 per cell (approximately 82 shelters per episode). While in a shelter, the agent takes no damage from enemies. Shelters do not protect against starvation.

Each episode uses a different random seed, producing a unique shelter layout and initial entity placement. Enemy and food spawning continues throughout the episode, maintaining population pressure.

\subsection{Observation Space}

The agent observes through an $15 \times 15$ window (radius 7) centered on the agent. The observation consists of three binary spatial channels:
\begin{enumerate}[nosep]
\item Enemy presence
\item Food presence
\item Shelter presence
\end{enumerate}
Cells outside the grid boundary are zero. Three scalar features are appended: $\frac{h}{h_{\max}}$, $\frac{u}{u_{\max}}$, and a binary shelter indicator. The full single-frame observation is a $678$-dimensional vector ($3 \times 15 \times 15 = 675$ spatial entries plus $3$ scalar features). The 4-frame stacked variant introduced as the observation-side ablation in Section~\ref{sec:ablation} concatenates the spatial channels of the four most recent frames, raising the spatial channel count from $3$ to $12$. The spatial component therefore grows by an exact factor of $4$ (from $675$ to $2{,}700$), but the three scalar features are taken from the current frame only, since they reflect the cumulative state of the agent and do not benefit from temporal stacking. Thus, the total observation dimension grows from $678$ to $2{,}703$, by a factor of $2703 / 678 \approx 3.99$. Figure~\ref{fig:obs_space} illustrates the layout.

\begin{figure}[h]
\centering
\includegraphics[width=\linewidth]{figures/observation_space}
\caption{Observation structure. The agent (blue dot at the window center) sees three binary $15 \times 15$ grids corresponding to enemy, food, and shelter presence within radius 7, plus three scalar features summarizing the agent's internal state. Section~\ref{sec:ablation} crosses this single-frame observation with a 4-frame stacked variant ($12$ spatial channels) as the observation-side ablation.}
\label{fig:obs_space}
\end{figure}

\subsection{Action Space}

The agent has five discrete actions: stay in place, or move one cell up, down, left, or right. Movement is clamped to the grid boundaries.

\subsection{Reward Structure}
\label{sec:reward_structure}

The reward at each step is a weighted sum of up to eleven components:
\begin{equation}
\label{eq:reward}
\begin{aligned}
r_t ={}& w_\text{food}\, i_t^{(\text{eat})}
    + w_\text{edm}\, d_t^{(\text{e})}
    + w_\text{sdm}\, d_t^{(\text{s})}
    + w_\text{death}\, i_t^{(\text{die})} \\
& + w_\text{hunger}\, i_t^{(\text{low})}
    + w_\text{hprop}(1 - u_t)
    + w_\text{tick} \\
& + w_\text{fprox}\, \Delta\phi_t^{(\text{f})}
    + w_\text{eprox}\, \Delta\phi_t^{(\text{e})}
    + w_\text{sprox}\, \Delta\phi_t^{(\text{s})} \\
& + w_\text{shelter}\, i_t^{(\text{safe})},
\end{aligned}
\end{equation}
where $w_{\cdot} \in \mathbb{R}$ are tunable scalar weights; $i_t^{(\text{eat})}, i_t^{(\text{die})}, i_t^{(\text{low})}, i_t^{(\text{safe})} \in \{0,1\}$ are binary indicators for food eaten, death, low hunger, and shelter safety; $d_t^{(\text{e})}$ and $d_t^{(\text{s})}$ are enemy damage and starvation damage; and $u_t \in [0, 1]$ is the hunger ratio.

The three proximity terms reward per-step changes in closeness to a target, rather than absolute closeness. For target type $x \in \{\text{f}, \text{e}, \text{s}\}$ (food, enemy, shelter), define the closeness score
\begin{equation}
\phi_t^{(x)} = \frac{d_\text{max} - d_t^{(x)} + 1}{d_\text{max} + 1} \in [0, 1],
\end{equation}
where $d_t^{(x)}$ is the Manhattan distance to the nearest visible entity of type $x$ and $d_\text{max} = 2r_\text{obs}$ (maximum Manhattan distance to any cell within the observation radius). If no entity is visible, $\phi_t^{(x)} = 0$. The per-step closeness change $\Delta\phi_t^{(x)} = \phi_t^{(x)} - \phi_{t-1}^{(x)}$ is positive when the agent moved closer to the target in that step and negative when it moved away. We reward this per-step change rather than the absolute closeness because rewarding absolute closeness would also reward standing still next to a target (for example, hovering next to a food tile without eating it). The per-step change form gives zero reward for staying in place, so it only pays the agent for active motion toward the target.

Two proximity terms can be conditionally gated on the hunger ratio. When conditional gating is enabled, the food proximity fires only when $u_t < u_\text{thresh}$ (hungry) and the shelter proximity fires only when $u_t \geq u_\text{thresh}$ (well-fed). This encodes a behavioral mode switch through the reward function: the agent receives a spatial gradient toward food when it needs to eat and toward shelter when it does not.

Table~\ref{tab:default_rewards} lists all reward components and their baseline weights.

\begin{table}[h]
\centering
\caption{Reward components and baseline weights. Components marked ``disabled'' are set to zero in the baseline and activated in specific experimental configurations.}
\label{tab:default_rewards}
\small
\begin{tabular}{llr}
\toprule
Component & Condition & Baseline \\
\midrule
Food eaten ($w_\text{food}$)          & Agent steps on food tile              & $+5.0$ \\
Enemy damage ($w_\text{edm}$)         & Per unit of damage received           & $-0.1$ \\
Starvation damage ($w_\text{sdm}$)    & Per unit of starvation damage         & $-1.0$ \\
Death ($w_\text{death}$)              & Health reaches zero                   & $-10.0$ \\
Low hunger ($w_\text{hunger}$)        & Per tick when $u_t < u_\text{thresh}$ & disabled \\
Hunger proportional ($w_\text{hprop}$) & Continuous, scales with $(1 - u_t)$  & $-0.3$ \\
Food proximity ($w_\text{fprox}$)     & Change in closeness to nearest food this step    & $+0.1$ \\
Enemy proximity ($w_\text{eprox}$)    & Change in closeness to nearest enemy this step   & disabled \\
Shelter proximity ($w_\text{sprox}$)  & Change in closeness to nearest shelter this step & disabled \\
Shelter safety ($w_\text{shelter}$)   & Agent in shelter                      & $+0.1$ \\
Survival tick ($w_\text{tick}$)       & Per step alive                        & disabled \\
\bottomrule
\end{tabular}
\end{table}

These components create competing pressures. Collecting food requires movement, which costs hunger and exposes the agent to enemies. Shelters neutralize enemy damage but provide no food. The hunger-proportional penalty grows every tick, creating urgency to eat that the agent cannot escape through passivity.

A subtle design consideration is the relative magnitude of continuous per-tick rewards versus sparse event rewards. A per-tick reward of $+0.1$ accumulated over a 700-step episode produces $+70$, which can dominate a food reward of $+5.0$. This creates incentives for passive survival: the agent learns to avoid risk rather than seek food. Similarly, a food proximity reward that is too large relative to the food-eaten reward incentivizes hovering near food without eating it, since eating removes the food and terminates the proximity stream. These failure modes motivate the experimental configurations in Section~\ref{sec:experiments}.

These weights are the primary independent variable in our experiments.

% =============================================================================
\section{Agent Architecture}
\label{sec:architecture}
% =============================================================================

\subsection{Network}
\label{sec:network}

Since the present work studies reward shape and observation-space changes rather than architecture, the architecture must be adequate for the task but does not need to be optimal. The key requirement is sufficient representational capacity so that the architecture is not the bottleneck. A convolutional encoder is the natural choice because the observation has spatial structure ($3k \times 15 \times 15$ grid for frame-stack size $k$). The Dueling head is inherited from the Rainbow-lite recipe. Channel counts and hidden dimensions follow standard values from the literature rather than task-specific tuning. For single-frame input ($k=1$), the network has $970{,}438$ trainable parameters for a $678$-dimensional input with $5$ output actions; the only architectural change in the frame-stacked cells ($k=4$) is the input depth of the first convolution. This provides ample capacity for the decision problem. Specific architectural choices (depth, width, kernel sizes) are not expected to materially affect the comparisons reported below.

The encoder splits the input into a $3 \times 15 \times 15$ spatial grid and 3 scalar features. Three convolutional layers process the spatial grid:
\begin{enumerate}[nosep]
\item $\text{Conv2d}(3 \to 32,\; 3 \times 3,\; \text{stride}\ 1,\; \text{pad}\ 1)$ \textrightarrow\ ReLU \textrightarrow\ LayerNorm
\item $\text{Conv2d}(32 \to 64,\; 3 \times 3,\; \text{stride}\ 1,\; \text{pad}\ 1)$ \textrightarrow\ ReLU \textrightarrow\ LayerNorm
\item $\text{Conv2d}(64 \to 64,\; 3 \times 3,\; \text{stride}\ 2,\; \text{pad}\ 0)$ \textrightarrow\ ReLU
\end{enumerate}
The third layer reduces the spatial dimensions from $15 \times 15$ to $7 \times 7$, producing a 3136-dimensional feature vector. This is concatenated with the 3 scalar features and projected through a fully connected layer ($3139 \to 256$) with ReLU and LayerNorm.

The Dueling head splits the 256-dimensional representation into two streams. The value stream maps $256 \to 128 \to 1$. The advantage stream maps $256 \to 128 \to 5$. Q-values are computed as $Q(s,a) = V(s) + A(s,a) - \bar{A}(s)$.

Convolutional layers use Kaiming uniform initialization. Fully connected layers use Xavier uniform initialization. The output layers of both Dueling streams use orthogonal initialization with gain 0.01 in order to produce near-zero initial Q-values and avoid large initial TD errors.

\subsection{Training Algorithm}
\label{sec:training_algorithm}

Algorithm~\ref{alg:training} summarizes the training procedure.

\begin{algorithm}[h]
\caption{Rainbow-lite DQN training.}
\label{alg:training}
\begin{algorithmic}[1]
\State Initialize policy network $\theta$, target network $\theta^- \leftarrow \theta$
\State Initialize prioritized replay buffer $\mathcal{D}$, $n$-step accumulator
\State Collect 10{,}000 transitions with random policy into $\mathcal{D}$
\For{$t = 1, \ldots, T$}
    \State Select $a_t$ via $\varepsilon$-greedy from $Q(\cdot; \theta)$, decay $\varepsilon$
    \State Execute $a_t$, observe $r_t, s_{t+1}, d_t$
    \State Compute $n$-step return; commit to $\mathcal{D}$ when ready
    \State Sample batch from $\mathcal{D}$ with PER; compute importance-sampling (IS) weights $w_i$
    \State $y_i = R_i^{(n)} + (1 - d_i)\,\gamma^n\, Q\!\left(s_i',\, \arg\max_{a'} Q(s_i', a'; \theta);\, \theta^-\right)$
    \State $L = \frac{1}{B}\sum_i w_i \cdot \text{Huber}(Q(s_i, a_i; \theta),\; y_i)$
    \State Update $\theta$ via AdamW with gradient clipping; step LR scheduler
    \State $\theta^- \leftarrow (1-\tau)\,\theta^- + \tau\,\theta$
    \State Update priorities in $\mathcal{D}$ from new TD errors
\EndFor
\end{algorithmic}
\end{algorithm}
Here, $d_i \in \{0,1\}$ indicates whether sampled transition $i$ is terminal, and $B$ denotes the batch size.
\paragraph{Experience collection.} At each step, the transition $(s_t, a_t, r_t, s_{t+1}, d_t)$ enters an $n$-step accumulator. Once $n = 3$ transitions have accumulated, the oldest is committed to the replay buffer with return $R_t^{(3)} = r_t + \gamma r_{t+1} + \gamma^2 r_{t+2}$, bootstrap state $s_{t+3}$, and discount factor $\gamma^3$. At episode termination (agent death), remaining entries are flushed with shortened returns. On episode truncation (time limit), pending entries are discarded to prevent cross-episode contamination of $n$-step returns.

The distinction between termination and truncation matters for the bootstrap target. A terminated episode (death) sets $d = 1$, zeroing the bootstrap term. A truncated episode sets $d = 0$, allowing bootstrapping from the final state since the agent was still alive.

\paragraph{Replay buffer.} Transitions are stored in pre-allocated numpy arrays. A SumTree data structure \citep{Schaul2016} supports efficient ($O(\log n)$) priority updates and proportional sampling. Stratified sampling divides $[0, P_\text{total}]$ into equal segments and draws one sample per segment, reducing variance compared to pure proportional sampling.

\paragraph{Optimization.} AdamW with decoupled weight decay ($10^{-4}$) on convolutional and linear weights; bias and normalization parameters are exempt. Gradients are clipped to norm 1.0. The learning rate is held constant.

\paragraph{Exploration.} $\varepsilon$-greedy with linear decay from 1.0 to 0.05 over the first half of training. The relatively high final $\varepsilon$ ensures continued exploration in the stochastic environment.

\subsection{Hyperparameters}
\label{sec:hyperparameters}

Table~\ref{tab:hyperparams} lists all hyperparameters. Values follow established defaults from the literature rather than environment-specific tuning. The frame-stack size $k$ is the observation-side ablation variable (Section~\ref{sec:ablation}); every other hyperparameter is held constant across all twenty-four ablation cells.

\begin{table}[h]
\centering
\caption{Agent hyperparameters.}
\label{tab:hyperparams}
\begin{tabular}{lll}
\toprule
Parameter & Value & Reference \\
\midrule
\multicolumn{3}{l}{\textit{Optimization}} \\
Learning rate              & $10^{-4}$ (constant) & \\
Weight decay               & $10^{-4}$         & \citet{Loshchilov2019} \\
Batch size                 & 128               & \\
Gradient clip norm         & 1.0               & \\
\midrule
\multicolumn{3}{l}{\textit{Replay buffer}} \\
Buffer capacity            & 250{,}000         & \\
Min.\ buffer size (warmup) & 10{,}000          & \citet{Mnih2015} \\
PER $\alpha$               & 0.6               & \citet{Schaul2016} \\
PER $\beta_0$              & 0.4               & \citet{Schaul2016} \\
$\beta$ annealing steps    & $T/2$             & \\
$n$-step returns           & 3                 & \citet{Hessel2018} \\
\midrule
\multicolumn{3}{l}{\textit{Target network}} \\
Polyak $\tau$              & 0.002             & \citet{Lillicrap2016} \\
\midrule
\multicolumn{3}{l}{\textit{Exploration}} \\
$\varepsilon$ start / end  & 1.0 / 0.01        & \citet{Mnih2015} \\
$\varepsilon$ decay steps  & 500{,}000         & \\
\midrule
\multicolumn{3}{l}{\textit{Network}} \\
CNN channels               & (32, 64, 64)      & \\
Merge hidden dim           & 256               & \\
Dueling head hidden dim    & 128               & \\
\midrule
\multicolumn{3}{l}{\textit{Observation (ablation variable)}} \\
Frame stack $k$            & $\{1, 4\}$        & \citet{Mnih2015} \\
\midrule
\multicolumn{3}{l}{\textit{Training}} \\
Discount $\gamma$          & 0.99              & \\
Total timesteps $T$        & 2{,}000{,}000     & \\
Eval frequency             & 5{,}000 steps     & \\
Eval episodes (in-training) & 10               & \\
Final benchmark episodes   & 100 (seeds 1000--1099) & \\
\bottomrule
\end{tabular}
\end{table}

% =============================================================================
\section{Experiments}
\label{sec:experiments}
% =============================================================================

\subsection{Experimental Setup}
\label{sec:setup}

Each ablation cell trains for $T = 2{,}000{,}000$ timesteps (a few hours on consumer GPUs). Episodes are capped at $1{,}000$ ticks (max\_steps in the world configuration). The cap is required by the protocol: an arbitrarily-foraging-and-fleeing policy could in principle survive indefinitely, which would prevent benchmark termination. We chose $1{,}000$ specifically because it leaves roughly $30\%$ headroom above the heuristic's mean survival of $\approx 768$ ticks, the comparison reference for this work. As an independent verification that this cap does not bias the headline comparison, we additionally train four extra seeds of \textsc{minimal\_fs} (the cell that comes closest to the heuristic) with the cap raised to $50{,}000$ ticks both during training and at benchmark; Section~\ref{sec:cap_robustness} reports the comparison. Greedy in-training evaluation runs every $5{,}000$ steps with $10$ episodes ($\varepsilon = 0$). Final benchmarks evaluate the best in-training checkpoint over $100$ episodes with fixed seeds ($1000$--$1099$). We report mean survival ticks (with $95\%$ confidence interval), mean food eaten, mean damage taken, and death rate. Each cell is run with four random seeds ($42$, $43$, $44$, $45$). The agent architecture, all hyperparameters in Table~\ref{tab:hyperparams}, and the training budget are held constant; the only variables are the reward configuration (Section~\ref{sec:reward_configs}) and the frame-stack size $k \in \{1, 4\}$ (Section~\ref{sec:ablation}).

\paragraph{Heuristic baseline.} A scripted agent provides a non-learned reference. The heuristic camps in the nearest shelter and leaves only when the hunger ratio drops below $0.5$, at which point it moves toward the nearest food. It flees from enemies within Manhattan distance $5$. Over $100$ evaluation episodes, this agent survives a mean of $768 \pm 19$ ticks ($95\%$ CI), eats $1.19$ food per episode, and dies in $94\%$ of episodes. The remaining $6\%$ reach the $1{,}000$-tick cap. The heuristic is rerun on the same evaluation seeds for every cell; $p$-values reported below are from two-sided Welch $t$-tests treating heuristic and DQN samples as independent.

\subsection{Failure Modes}
\label{sec:failure_modes}

Early experiments revealed three degenerate policies that emerged reliably under certain reward configurations. These modes are not mutually exclusive: all three appeared within a single baseline training run at different stages and in different episodes, with the agent alternating between hovering near food, seeking enemies, and oscillating in place depending on the local state distribution.

\paragraph{Oscillation.} The agent remains passive or oscillates between adjacent cells. This behavior was observed under several reward configurations and its precise cause requires further investigation.

\paragraph{Food hovering.} The agent moves near food but does not eat it. This occurs when the food proximity reward exceeds the hunger penalty. An agent 5 cells from food with $w_\text{prox} = 0.5$ earns $+0.33$ per tick from proximity, while the average hunger penalty is $-0.15$ per tick, yielding a positive net reward for inaction. Eating the food would remove the proximity stream. Heuristic H2 addresses this by bounding $w_\text{prox}$ below $|w_\text{hprop}|/2$.

\paragraph{Suicide.} When the proximity gradient is too weak to guide the agent toward food, the agent cannot learn to forage. Without foraging, every tick accumulates hunger-proportional penalty. A stationary agent surviving 700 ticks until starvation accumulates approximately $-135$ from $w_\text{hprop}$ and $-100$ from $w_\text{sdm}$, totaling $-235$ before the death penalty. The agent discovers that dying quickly is preferable: seeking enemies yields a fast death at total cost $\approx -24$ (a few enemy hits plus the death penalty), roughly 10$\times$ less than slow starvation. This is a rational response to a reward function where living is more expensive than dying.

The food hovering and suicide failure modes are clear instances of reward hacking~\citep{Amodei2016}: the agent optimizes the reward function as specified, but the resulting behavior diverges from the designer's intent. The cause of oscillation is less certain and requires further investigation. These failure modes motivate the design heuristics derived in Section~\ref{sec:reward_configs}; Section~\ref{sec:results} quantifies how much of each cell's evaluation distribution falls into each failure mode.

\subsection{Reward Configurations}
\label{sec:reward_configs}

The failure modes above show that poorly chosen reward weights produce degenerate policies. Based on the per-step reward dynamics, we derive four design heuristics (H1--H4) that help avoid the food-hovering and suicide failure modes. These compare quantities at the same timestep, so discounting does not affect them.

Let $D = 2 r_\text{obs}$ denote the observation diameter, $c$ the per-step movement hunger cost, $d_e$ the enemy contact damage, and $h_s$ the per-tick starvation damage.
\begin{enumerate}[label=\textbf{H\arabic*},nosep,leftmargin=2.5em]
  \item \emph{Anti-stasis.} One step toward food must yield more proximity reward than the hunger penalty from the movement cost:
    $w_\text{prox} > |w_\text{hprop}| \cdot c \cdot (D+1) \,/\, u_\text{max}$.
  \item \emph{Anti-hovering.} Maximum proximity reward must not exceed the average hunger penalty:
    $w_\text{prox} < |w_\text{hprop}| / 2$.
  \item \emph{Enemy flee.} Damage penalty from one hit must exceed the hunger penalty from one evasive step:
    $|w_\text{edm}| \cdot d_e > |w_\text{hprop}| \cdot c \,/\, u_\text{max}$.
  \item \emph{Starvation escalation.} Per-tick penalty when starving must exceed the max hunger-proportional penalty:
    $|w_\text{sdm}| \cdot h_s > |w_\text{hprop}|$.
\end{enumerate}

The food-hovering failure mode maps to an H2 violation and the suicide failure mode to an H1 violation.

The reward-side ablation uses three named configurations. Two satisfy or violate specific design heuristics; one is a minimal-component reward that renders the heuristics vacuous by construction. We additionally include a single-seed illustration that demonstrates the empirical signature of an H1 violation. Table~\ref{tab:reward_configs} summarises them.

\begin{table}[h]
\centering
\caption{Reward configurations used in the experiments. Only weights that differ from Table~\ref{tab:default_rewards} are shown. \textsc{minimal} sets $w_\text{hprop} = 0$, which renders H1--H4 (all of which involve $w_\text{hprop}$) vacuous as constraints.}
\label{tab:reward_configs}
\small
\begin{tabular}{lll}
\toprule
Configuration & Modification & Role \\
\midrule
\textsc{baseline}            & (none --- default weights)                        & ablation cell, satisfies H1--H4 \\
\textsc{absolute\_proximity} & $\Delta\phi$ replaced by $\phi$ (no closeness change) & ablation cell, violates H2 by construction \\
\textsc{minimal}      & minimal four-component reward (Table~\ref{tab:minimal_config}) & ablation cell, heuristics-vacuous ($w_\text{hprop}=0$) \\
\textsc{weak\_proximity}     & $w_\text{prox} = 0.02$                            & illustration of an H1 violation (single seed) \\
\bottomrule
\end{tabular}
\end{table}

\paragraph{Baseline.} The default weights (Table~\ref{tab:default_rewards}) satisfy all four heuristics. The hunger-proportional penalty ($-0.3$) creates urgency, the food proximity gradient ($+0.1$) provides a directional signal, and the food, damage, shelter, and death components encode the rest of the survival objective. The configuration is dense and compositional: every entity type contributes a reward signal and every internal state has a corresponding term. This makes \textsc{baseline} the most informative reward in principle, but also the most exposed to interactions between components.

\paragraph{Absolute proximity.} The proximity terms are switched from per-step closeness change $\Delta\phi$ to absolute closeness $\phi$. Every other weight matches \textsc{baseline}. By construction this violates H2: an agent standing one cell from a food tile receives the maximum proximity reward indefinitely, while the average hunger penalty over the same interval is finite and smaller. This configuration directly tests the importance of the closeness-change reward shape introduced in Equation~\ref{eq:reward}: with absolute proximity, a hovering policy should be incentivised even though all other components are unchanged.

\paragraph{Minimal.} A minimal reward set with only three closeness-change proximity gradients and the death penalty:
\begin{itemize}[nosep,leftmargin=1.6em]
  \item food-proximity ($w_\text{fprox}=+0.15$), gated to fire only when $u_t < 0.5$ (hungry);
  \item shelter-proximity ($w_\text{sprox}=+0.15$), gated to fire only when $u_t \geq 0.5$ (well-fed);
  \item enemy-proximity ($w_\text{eprox}=-0.5$), always active;
  \item death ($w_\text{death}=-10$).
\end{itemize}
Every other weight is set to zero. The motivation is to cleanly separate the spatial gradients from the event-driven rewards (food eaten, damage taken, shelter safety) and from the always-on hunger penalty: each of the remaining three signals corresponds to a single objective, and the gating mirrors the heuristic's mode switch (forage when hungry, shelter when not). Removing $w_\text{hprop}$ also resolves a structural incompatibility in the baseline. With $\gamma = 0.99$ the present value of the per-step hunger-cost penalty over the remaining episode is approximately $0.15$, while the maximum food-proximity reward from one step is approximately $0.0067$; the movement penalty exceeds the approach reward by a factor of about $22$, so not moving is locally Q-value optimal regardless of hunger level. \textsc{minimal} removes $w_\text{hprop}$ entirely, eliminating this stasis equilibrium. Table~\ref{tab:minimal_config} lists the resulting weights.

\begin{table}[h]
\centering
\caption{\textsc{minimal} weights. All unlisted components from Equation~\ref{eq:reward} are set to zero. Food and shelter proximity are conditionally gated on the hunger threshold $u_\text{thresh} = 0.5$.}
\label{tab:minimal_config}
\small
\begin{tabular}{llr}
\toprule
Component & Condition & Weight \\
\midrule
Food proximity ($w_\text{fprox}$)    & Closeness change, active only when $u_t < 0.5$    & $+0.15$ \\
Shelter proximity ($w_\text{sprox}$) & Closeness change, active only when $u_t \geq 0.5$ & $+0.15$ \\
Enemy proximity ($w_\text{eprox}$)   & Closeness change, always active                   & $-0.50$ \\
Death ($w_\text{death}$)             & Terminal                                          & $-10.0$ \\
\bottomrule
\end{tabular}
\end{table}

\paragraph{Weak proximity (illustration).} The food proximity weight is reduced from $0.1$ to $0.02$, putting it below the H1 anti-stasis threshold ($|w_\text{hprop}| \cdot c \cdot (D+1) / u_\text{max} = 0.0225$). All other weights match \textsc{baseline}. This is included as a single-seed illustration of the suicide failure mode predicted by an H1 violation, not as part of the main ablation. Section~\ref{sec:weak_prox} reports the resulting agent behaviour qualitatively.

\subsection{Weak Proximity: an H1 Illustration}
\label{sec:weak_prox}

The H1 anti-stasis bound predicts that food proximity rewards below the per-step threshold should fail to produce viable foraging. We trained a single seed of \textsc{weak\_proximity} ($w_\text{fprox}=0.02$, all other weights matching \textsc{baseline}) for two million steps and benchmarked the result over $100$ episodes. The agent reached mean survival $647 \pm 33$ ticks with $0.66$ food per episode and a $99\%$ death rate. These aggregate statistics are very close to the same architecture trained on \textsc{baseline} ($651 \pm 41$ ticks, $0.70$ food, $98\%$ death rate). We include the run as an illustration that H1 is a per-step lower bound on the proximity weight, not a predictor of large aggregate effects, and as a reminder that the heuristics in Section~\ref{sec:reward_configs} are necessary conditions rather than predictions of policy quality. We do not draw conclusions from this single seed beyond noting that the H1 violation does not produce a markedly different aggregate outcome from \textsc{baseline} on this single training run.

\subsection{Reward $\times$ Observation Ablation}
\label{sec:ablation}

The main experiment crosses three reward configurations (\textsc{baseline}, \textsc{absolute\_proximity}, \textsc{minimal}) with two observation types (single-frame, $k = 1$, and 4-frame stacked, $k = 4$), four random seeds per cell. Each cell trains for two million timesteps under the hyperparameters in Table~\ref{tab:hyperparams}; the only varying parameters are the reward weights (Section~\ref{sec:reward_configs}) and $k$. The benchmark protocol is described in Section~\ref{sec:setup}.

\paragraph{Hypotheses.} This ablation tests the following hypotheses:
\begin{enumerate}[label=(H\arabic*),nosep,leftmargin=2.8em]
  \item \textsc{minimal} with frame stacking is the highest-performing of the six cells.
  \item \textsc{minimal} with frame stacking exceeds the heuristic.
  \item Frame stacking improves every reward configuration over its single-frame counterpart.
\end{enumerate}
The reasoning behind H1 and H2 was that the cleanest reward (\textsc{minimal}) combined with the temporal context required to infer enemy velocity (4-frame stacking) should produce the strongest policy in the experiment, and that this strongest policy would be sufficient to beat the scripted heuristic. H3 captured the prior expectation that adding observation history is at worst neutral: extra information that an RL agent can in principle ignore should not degrade learned performance.

\paragraph{Frame stacking.} 4-frame stacking concatenates the spatial channels from the last four time steps so the CNN can infer inter-frame position changes from which velocity can be derived. The base observation has $3$ spatial channels on a $15 \times 15$ grid; a 4-frame stack provides $12$. Scalar features (health, hunger, shelter indicator) are taken from the current step only, since they reflect the cumulative state of the agent. The $k=4$ observation has dimension $|O_4| = 3 \cdot 4 \cdot 225 + 3 = 2703$, an increase of $3.99\times$ over $|O_1|$.

\subsection{Results}
\label{sec:results}

Table~\ref{tab:cell_results} reports per-cell averages across four seeds. Table~\ref{tab:per_seed} reports the per-seed numbers. Figure~\ref{fig:ablation_grid} summarises survival, foraging activity, and the frame-stacking effect.

\begin{figure}[t]
\centering
\includegraphics[width=\linewidth]{figures/ablation_grid}
\caption{Reward $\times$ observation ablation. Each bar in the first two panels is the mean over four seeds; black dots show individual seeds; error bars are the standard error of the mean. Dashed lines mark the heuristic baseline ($768$ ticks survival, $1.19$ food per episode). The third panel shows the mean change in survival between the single-frame and 4-frame configurations of each reward. Frame stacking helps \textsc{minimal} substantially ($+71$) but hurts \textsc{baseline} ($-91$) and \textsc{absolute\_proximity} ($-42$).}
\label{fig:ablation_grid}
\end{figure}

\begin{table}[h]
\centering
\caption{Reward $\times$ observation ablation. Each cell aggregates four seeds ($42$, $43$, $44$, $45$). ``Diff'' is the mean across seeds of the per-seed (DQN $-$ heuristic) survival difference. ``Sig.\ losses'' counts the seeds (out of four) for which the heuristic significantly outperformed DQN at $p < 0.05$ on a two-sided Welch $t$-test against the matched-seed heuristic baseline. \textsc{weak\_proximity} is a single-seed reference and is reported separately. Numbers below \textsc{minimal\_fs} (cell with non-trivial cap-rate) are corroborated by an independent four-seed cap-$50{,}000$ training in Section~\ref{sec:cap_robustness}.}
\label{tab:cell_results}
\footnotesize
\setlength{\tabcolsep}{4pt}
\begin{tabular}{lrrrrr}
\toprule
Cell & Survival (mean) & Food & Death rate & Diff vs heuristic & Sig.\ losses \\
\midrule
Heuristic                          & $768 \pm 19$       & $1.19$ & $94\%$  & $0$       & --- \\
\midrule
\multicolumn{6}{l}{\textit{Single-frame ($k=1$)}} \\
\textsc{baseline}                  & $647$              & $0.69$ & $98\%$  & $-122$    & $4/4$ \\
\textsc{absolute\_proximity}       & $634$              & $0.39$ & $99\%$  & $-135$    & $4/4$ \\
\textsc{minimal}                   & $674$              & $0.09$ & $100\%$ & $-94$     & $4/4$ \\
\midrule
\multicolumn{6}{l}{\textit{4-frame stacked ($k=4$)}} \\
\textsc{baseline\_fs}              & $556$              & $4.86$ & $96\%$  & $-212$    & $4/4$ \\
\textsc{absolute\_proximity\_fs}   & $592$              & $4.75$ & $97\%$  & $-177$    & $4/4$ \\
\textsc{minimal\_fs}               & $\mathbf{745}$     & $2.14$ & $94\%$  & $\mathbf{-23}$ & $1/4$ \\
\midrule
\textsc{weak\_proximity} (1 seed)  & $647$              & $0.66$ & $99\%$  & $-122$    & $1/1$ \\
\bottomrule
\end{tabular}
\end{table}

\begin{table}[h]
\centering
\caption{Per-seed DQN survival. Heuristic survival is $768 \pm 19$ ticks. Each $p$ is a two-sided Welch $t$-test against the heuristic for that seed (italics indicate statistical ties at $p \geq 0.05$). Seed $45$ was trained later than $42$--$44$ as a confirmatory expansion of the original $3 \times 2 \times 3$ ablation.}
\label{tab:per_seed}
\footnotesize
\begin{tabular}{lrrrrl}
\toprule
Cell                              & seed 42 & seed 43 & seed 44 & seed 45 & per-seed verdict \\
\midrule
\textsc{baseline}                 & $623$ & $697$ & $632$ & $636$ & all heur wins ($p \leq 0.0002$) \\
\textsc{baseline\_fs}             & $551$ & $541$ & $552$ & $582$ & all heur wins ($p \leq 0.0001$) \\
\textsc{absolute\_proximity}      & $623$ & $651$ & $646$ & $617$ & all heur wins ($p \leq 0.0001$) \\
\textsc{absolute\_proximity\_fs}  & $600$ & $583$ & $590$ & $595$ & all heur wins ($p \leq 0.0001$) \\
\textsc{minimal}                  & $677$ & $670$ & $664$ & $685$ & all heur wins ($p \leq 0.0001$) \\
\textsc{minimal\_fs}              & \emph{$757$} & $731$ & \emph{$745$} & \emph{$746$} & $3/4$ tie ($p\!\in\!\{0.47,\ 0.11,\ 0.10\}$); seed $43$ heur wins ($p\!=\!0.009$) \\
\bottomrule
\end{tabular}
\end{table}

\paragraph{Failure-mode analysis.}

DESCRIPTION HERE

\begin{figure}[t]
\centering
\includegraphics[width=\linewidth]{figures/failure_modes}
\caption{(Left) failure-mode composition of the $100$-episode benchmark, averaged across the four seeds for each cell. (Right) mean shelter occupancy (fraction of survived ticks spent inside a shelter). The heuristic spends $\approx 82\%$ of its life in shelters; every DQN cell spends $<2\%$ regardless of reward configuration.}
\label{fig:failure_modes}
\end{figure}

\begin{table}[h]
\centering
\caption{Per-cell failure-mode fractions and shelter occupancy, in percent. Categories are mutually exclusive and sum to $100\%$ within rounding. ``Sh.\ occ.'' is the mean fraction of survived ticks spent inside a shelter, expressed as a percentage. Each row averages four seeds (the heuristic row is a single $100$-episode run).}
\label{tab:failure_modes}
\footnotesize
\setlength{\tabcolsep}{4pt}
\begin{tabular}{lrrrrrr}
\toprule
Cell & Cap & Enemy & Forage/die & Sheltered & Exposed & Sh.\ occ. \\
     &     & suicide & early       & passivity & stasis  &   \\
\midrule
Heuristic                          & $6.0$ & $0.0$  & $26.0$ & $\mathbf{57.0}$ & $11.0$ & $82.1$ \\
\midrule
\multicolumn{7}{l}{\textit{Single-frame ($k=1$)}} \\
\textsc{baseline}                  & $1.5$ & $6.3$  & $13.0$ & $0.5$           & $\mathbf{78.8}$ & $0.9$  \\
\textsc{absolute\_proximity}       & $0.5$ & $8.0$  & $8.3$  & $0.8$           & $\mathbf{82.5}$ & $1.0$  \\
\textsc{minimal}                   & $0.0$ & $9.5$  & $0.8$  & $0.8$           & $\mathbf{89.0}$ & $0.7$  \\
\midrule
\multicolumn{7}{l}{\textit{4-frame stacked ($k=4$)}} \\
\textsc{baseline\_fs}              & $2.8$ & $7.0$  & $\mathbf{76.8}$ & $0.0$  & $13.5$          & $0.6$  \\
\textsc{absolute\_proximity\_fs}   & $4.0$ & $6.5$  & $\mathbf{61.5}$ & $1.0$  & $27.0$          & $1.7$  \\
\textsc{minimal\_fs}               & $\mathbf{4.3}$ & $0.3$ & $\mathbf{52.0}$ & $0.0$ & $43.5$ & $0.8$  \\
\bottomrule
\end{tabular}
\end{table}

CONCLUSION HERE

Shelter occupancy is essentially zero across all six DQN cells (median $0.8\%$ of survived ticks; max $1.7\%$ for \textsc{absolute\_proximity\_fs}), confirming that no cell has learned to use the shelters as the heuristic does.

\paragraph{Learning trajectories.}
FIX FIGURE PLACEMENT
Figure~\ref{fig:learning_curves} plots the in-training $10$-episode evaluation survival as a function of training step for each of the six ablation cells. The trajectories support three observations relevant to the experimental design. First, every cell reaches its operating range well before the two-million-step budget elapses: each cell's curve flattens by step $1.0$--$1.4$M, and the second half of training contributes drift rather than a discernible upward trend in any cell. Second, the trajectories are tightly clustered across seeds within a cell, which corroborates the per-seed homogeneity reported in Table~\ref{tab:per_seed} and indicates that the variance across cells is dominated by the (reward, observation) pair rather than by random initialisation. Third, the displayed values use a $10$-eval rolling mean over $10$ episodes per eval, so they are noisier than the $100$-episode benchmark numbers reported in Table~\ref{tab:cell_results}; the benchmark targets \texttt{model\_best.pt}, which was selected from the highest in-training eval and tends to lie above the smoothed late-training operating range (most visibly for \textsc{minimal\_fs}, whose smoothed late curve sits near $620$--$660$ ticks while \texttt{model\_best.pt} benchmarks at $744$). We did not observe a plasticity-collapse signature \citep{Lyle2023, Dohare2024} in any cell: the curves saturate but do not decay over the course of training. We take the absence of an upward trend in the final third of training as the operational justification for terminating at $T = 2{,}000{,}000$ steps; a longer budget would be necessary only to test the fs-leaky-reward hypothesis discussed in Section~\ref{sec:discussion}.

\begin{figure}[t]
\centering
\includegraphics[width=\linewidth]{figures/learning_curves}
\caption{In-training evaluation survival vs training step for the six ablation cells. Each line is one seed, smoothed with a $10$-eval rolling mean ($50{,}000$-step window). The dashed grey line at $768$ ticks marks the heuristic baseline. All six cells have flattened by approximately $1.0$--$1.4$M steps; the final third of training does not produce an upward trend in any cell, so we terminate training at $T = 2{,}000{,}000$ steps.}
\label{fig:learning_curves}
\end{figure}

\subsection{Summary of Observations}
\label{sec:analysis}

We organise the analysis around the three hypotheses.

\paragraph{(H1) \textsc{minimal\_fs} is the highest-performing cell. HELD.}
Mean survival across the six cells ranges from $556$ (\textsc{baseline\_fs}) to $745$ (\textsc{minimal\_fs}). \textsc{minimal\_fs} is the only cell that comes close to the heuristic; the second-best cell is \textsc{minimal} (single-frame) at $674$, trailing by $71$ ticks. The within-cell standard error of \textsc{minimal\_fs} is small enough that no other cell is consistent with matching its mean.

\paragraph{(H2) \textsc{minimal\_fs} exceeds the heuristic. NOT HELD.}
Frame stacking lifts \textsc{minimal} from a mean of $674$ to $745$ ticks, a $+71$-tick gain. Three of the four frame-stacked seeds are statistically tied with the heuristic ($p = 0.47$, $0.11$, $0.10$); the fourth significantly loses ($p = 0.009$). The mean survival difference of $-23$ ticks is consistent with parity, not superiority. We had expected clear superiority; we observed parity.

\paragraph{(H3) Frame stacking improves every reward configuration. NOT HELD.}
With $\Delta = \text{mean}(\text{fs}) - \text{mean}(\text{sf})$: $\Delta = +71$ for \textsc{minimal}, $\Delta = -42$ for \textsc{absolute\_proximity}, $\Delta = -91$ for \textsc{baseline} (Figure~\ref{fig:ablation_grid}, third panel). Frame stacking helps the cleanest reward and hurts both leakier ones. The direction reverses cleanly with reward shape; this is the most striking departure from our prior expectation. We discuss plausible mechanisms in Section~\ref{sec:discussion}.

\paragraph{Single-frame cells cluster well below the heuristic.}
Although not one of the three hypotheses, the single-frame results are worth reporting on their own. Mean survival across the three single-frame reward configurations is $674$ for \textsc{minimal}, $647$ for \textsc{baseline}, and $634$ for \textsc{absolute\_proximity}, against the heuristic's $768$. Every single-frame seed is significantly outperformed by the heuristic at $p < 0.001$. The reward configuration matters at the margin (range $\approx 40$ ticks across the three rewards) but not enough to bring DQN within reach of the heuristic without temporal context. \textsc{baseline\_fs} and \textsc{absolute\_proximity\_fs} are also significantly below their single-frame counterparts: \textsc{baseline\_fs} is $91$ ticks below \textsc{baseline} ($556$ vs $647$); \textsc{absolute\_proximity\_fs} is $42$ ticks below \textsc{absolute\_proximity} ($592$ vs $634$). Frame stacking does not merely fail to rescue these reward configurations, but actively makes them worse.

\paragraph{Foraging over Safety.}
Frame stacking changes the foraging rate substantially: from $0.09$ to $2.14$ for \textsc{minimal}, from $0.39$ to $4.75$ for \textsc{absolute\_proximity}, and from $0.69$ to $4.86$ for \textsc{baseline}. The frame-stacked reward gradients teach the agent to seek food more aggressively, but in the leaky reward configurations (\textsc{baseline}, \textsc{absolute\_proximity}) the resulting policy spends more ticks exposed to enemies and dies sooner. Food-eating behaviour is unlocked; safety behaviour is not preserved. Only \textsc{minimal\_fs} reaches a foraging rate consistent with sustained survival ($2.14$ food per episode at survival within $23$ ticks of the heuristic, against the heuristic's $1.19$ food at $768$ ticks).

\paragraph{Shelters are not used.}
SHELTERS NOT USED

\paragraph{No DQN cell beats the heuristic.}
Across the twenty-four ablation-cell seeds plus the single-seed illustration and the four cap-$50{,}000$ robustness seeds, no $100$-episode benchmark significantly outperforms the heuristic. The closest cap-$1{,}000$ seeds are three of the four frame-stacked \textsc{minimal} seeds, which tie the heuristic at $p = 0.47$, $p = 0.11$, and $p = 0.10$; the fourth (seed $43$) significantly loses at $p = 0.009$.

\paragraph{Closeness-change vs absolute closeness.}
The contrast between \textsc{baseline} and \textsc{absolute\_proximity} isolates the effect of the closeness-change reward shape. Holding all other weights constant and switching the proximity term from $\Delta\phi$ to $\phi$ degrades survival by $13$ ticks at single-frame and by $35$ ticks at frame-stacked. The absolute-closeness form pays the agent for being near a target (including for standing still next to it), and the resulting policies eat more food but die earlier. The closeness-change form, used in \textsc{minimal}, is the only reward shape under which frame stacking improves survival.

\paragraph{Episode-cap robustness check.}
\label{sec:cap_robustness}
The $1{,}000$-tick cap is the only protocol parameter chosen ad hoc rather than from the literature, so we verify directly that it is not biasing the headline comparison. The heuristic reaches the cap in $6\%$ of episodes; \textsc{minimal\_fs} hits it in $1$--$10\%$ of episodes per seed (mean $5\%$); every other DQN cell hits it in at most $3\%$ of episodes and is far enough below the heuristic that small cap-induced bias cannot affect its conclusion. To verify that the parity-not-superiority verdict for \textsc{minimal\_fs} is not an artefact of training under truncation, we trained four additional seeds (seeds $45$--$48$) of the same configuration with the cap raised from $1{,}000$ to $50{,}000$ ticks both during training and at benchmark, holding every other hyperparameter fixed. The same $100$-episode evaluation seeds ($1000$--$1099$) are used. The heuristic is also re-benchmarked at this cap to keep the comparison apples-to-apples.

Three observations follow. First, the longest individual episode across all four DQN seeds and the heuristic is $1{,}516$ ticks (a single \textsc{minimal\_fs}-cap50k seed-$48$ episode), and no policy approaches an order of magnitude of the lifted cap. Second, the heuristic's mean survival is essentially unchanged ($768.8 \pm 19.1$ at cap $1{,}000$ vs $769.9 \pm 19.9$ at cap $50{,}000$), confirming that the heuristic's small ($\approx 6\%$) cap-rate at $1{,}000$ ticks did not materially bias its mean. Third, the four cap-$50{,}000$ seeds reach mean survival $735$--$791$ ticks (per-seed numbers in Table~\ref{tab:cap_robustness}); three of the four are statistically tied with the heuristic at $p \geq 0.20$, and the fourth significantly loses ($p = 0.03$). Averaging across the four seeds gives a mean of $767.5$ ticks against the heuristic's $769.9$, qualitatively the same parity-not-superiority pattern as the four cap-$1{,}000$ seeds (mean $745$, $3/4$ ties). The mean shift between cap-$1{,}000$ and cap-$50{,}000$ training is $\approx +23$ ticks, comparable to the cell's per-seed standard deviation; the relative ranking against the heuristic is unchanged.

\begin{table}[h]
\centering
\caption{Cap-$50{,}000$ robustness benchmark for \textsc{minimal\_fs}. Heuristic re-benchmarked at the same cap. ``Diff'' is DQN $-$ heuristic; $p$ from a two-sided Welch $t$-test. \textsc{minimal\_fs}-cap50k seeds were trained for $\approx 1.7$M steps; the in-training eval had plateaued by then and \texttt{model\_best.pt} was already selected from earlier-step peaks (see Section~\ref{sec:results} learning trajectories paragraph for the analogous shape on the cap-$1{,}000$ runs).}
\label{tab:cap_robustness}
\footnotesize
\begin{tabular}{lrrrrl}
\toprule
Cell                                              & Survival           & Diff   & $p$    & Max ep.\ & Verdict \\
\midrule
Heuristic (cap $50{,}000$)                        & $769.9 \pm 19.9$   & ---    & ---    & $1{,}072$ & ---   \\
\textsc{minimal\_fs}-cap50k seed $45$             & $779.3 \pm 29.2$   & $+9.3$ & $0.60$ & $1{,}337$ & tie   \\
\textsc{minimal\_fs}-cap50k seed $46$             & $763.8 \pm 29.2$   & $-6.2$ & $0.73$ & $1{,}153$ & tie   \\
\textsc{minimal\_fs}-cap50k seed $47$             & $791.4 \pm 27.5$   & $+21.4$& $0.22$ & $1{,}220$ & tie   \\
\textsc{minimal\_fs}-cap50k seed $48$             & $735.4 \pm 24.7$   & $-34.6$& $0.03$ & $1{,}516$ & heur wins \\
\bottomrule
\end{tabular}
\end{table}

The remaining cells are evaluated only at the $1{,}000$-tick cap; their gap to the heuristic is large enough that any plausible cap-bias would not change their conclusions.

% =============================================================================
\section{Discussion}
\label{sec:discussion}
% =============================================================================

\paragraph{The hypothesis we set out to test.}
The thesis was framed around the hypothesis that frame stacking applied to a clean minimal reward (\textsc{minimal}) would lift the agent above the heuristic. The data partially support this. \textsc{minimal\_fs} is the best-performing cell, by a clear margin over every other reward $\times$ observation combination, and frame stacking does add $+71$ ticks of mean survival on this clean reward. But \textsc{minimal\_fs} reaches only statistical parity with the heuristic in three of four seeds and significantly loses in the fourth (mean diff $-23$ ticks). The qualitative reading of the original hypothesis survives; the quantitative claim that the agent would beat the heuristic does not.

\paragraph{Frame stacking's reward-dependent direction.}
The most surprising result is that frame stacking degrades survival in two of three reward configurations. With the cleanest reward (\textsc{minimal}) it helps by $71$ ticks on average; with the leaky reward (\textsc{absolute\_proximity}) it hurts by $42$ ticks; with \textsc{baseline} it hurts by $91$ ticks. In principle, additional observation channels should weakly improve the agent: extra information can be ignored if it is not useful. The data suggests this is not what happens here. HYPOTHESES HERE

\paragraph{Limitations.}
\begin{itemize}[nosep,leftmargin=1.6em]
  \item \emph{Single environment.} Our findings are only observed in this specific environment and a single task. To generalize, replication on additional environments would be necessary.
  \item \emph{Architecture held fixed.} We do not run a network-architecture sweep. Capacity-matched comparisons between $k=1$ and $k=4$ encoders are a natural follow-up; without them we cannot rule out that the degradation under leaky rewards is purely an artefact of holding the post-conv architecture fixed across $k$.
  \item \emph{Training budget.} Due to limited computational budget, each cell was only trained for two million steps. Longer training could test whether certain configurations achieve better performance.
  \item \emph{Single seed for the H1 illustration.} Due to \textsc{weak\_proximity} being trained on only one seed, we do not draw conclusions from it beyond the qualitative note in Section~\ref{sec:weak_prox}.
\end{itemize}

% =============================================================================
\section{Conclusion}
\label{sec:conclusion}
% =============================================================================

We ran a $3 \times 2 \times 4$ ablation that crossed three reward configurations (\textsc{baseline}, \textsc{absolute\_proximity}, \textsc{minimal}) with two observation types (single-frame and 4-frame stacked), four random seeds per cell. We added a single-seed illustrative run of \textsc{weak\_proximity} for the H1 anti-stasis bound. All training was pure RL with no demonstrations.

\paragraph{What we hypothesised, what we observed.}
We hypothesised three things: that \textsc{minimal\_fs} would be the best-performing cell (H1); that it would exceed the heuristic (H2); and that frame stacking would improve every reward configuration over its single-frame counterpart (H3). H1 held: \textsc{minimal\_fs} reaches mean survival $745$ ticks, the highest of any cell. H2 did not hold: \textsc{minimal\_fs} reached statistical parity with the heuristic in three of four seeds and significantly lost in the fourth, but did not exceed it. H3 also did not hold: frame stacking improves \textsc{minimal} ($+71$ ticks) but degrades \textsc{absolute\_proximity} ($-42$) and \textsc{baseline} ($-91$). The qualitative direction of our prior expectation survives --- clean reward and temporal context together produce the best policy --- but the strongest claims about magnitude and universality do not.

\paragraph{The reward-dependent direction.}
Beyond falsifying H3, the data reveal a directional structure to frame stacking's effect: it improves the cleanest reward configuration and degrades both leakier ones, by margins much larger than within-cell variance. Adding observation history is therefore not an unconditional improvement on this task; under leaky reward designs it actively hurts the resulting agent. Plausible mechanisms are discussed in Section~\ref{sec:discussion}.

\paragraph{Position of the contribution.}
This work is a single-environment characterisation of the reward $\times$ observation interaction in DQN, not a positive ``DQN beats heuristic'' result. The contribution is: \begin{itemize}
  \item a derivation of four reward-design heuristics (H1--H4) from environment mechanics
  \item a $3 \times 2 \times 4$ reward $\times$ observation $\times$ seed ablation that tests the three hypotheses and reveals the reward-dependent direction of frame stacking
\end{itemize}

\paragraph{Future work.}
Four directions follow naturally from these results. First, replicate the ablation on a second multi-objective environment to test the generality of the reward-dependent direction. Second, run capacity-matched architecture comparisons between $k=1$ and $k=4$ to isolate whether the degradation under leaky rewards is an architectural effect (capacity dilution) or a reward effect (more surface for hacking). Third, train the frame-stacked leaky-reward cells for substantially longer than two million steps to test the training-budget hypothesis.

% =============================================================================
% References
% =============================================================================
\newpage
\bibliographystyle{plainnat}
\bibliography{refs}

\end{document}
